% Unit 14 — the factorial reduction
\input{technique-note.tex}
\title{Integration Bee --- Unit 14:\\ A factorial reduction}
\author{UK Uni Integration Bee prep\\ The integer Gamma integral, then Yang \S2.4}
\date{}
\begin{document}
\maketitle
\section*{Status}
Tags: \texttt{yang:2.4}, \texttt{gamma}; optional syllabus nodes \texttt{beta} and \texttt{euler-mascheroni}.
The identity $n!=\int_0^{\infty}x^n e^{-x}\,dx$ is the integer case of the Gamma function. Yang \S2.4 is one writeup. The reduction itself is parts: the index drops by 1. Unit 20 names the same family and stores neither $J_2$ nor $J_4$.
The optional recognition lessons below host Beta and Euler's constant. They add no Round 1 gate or machine drill id. A ready lesson does not assert a learner's mastery.
\section{Official anchors}
\begin{itemize}
  \item Parts on $[0,\infty)$, any calculus text. Unit 20 for the trigonometric $I_n$.
  \item Syllabus §2 Gamma, as the name of this integral.
  \item Yang Topics Guide 2nd Ed §2.4 Gamma Function Integral.
  \item Syllabus §2 Beta and Euler--Mascheroni; coach cards \path{docs/techniques/concepts/beta.md} and \path{docs/techniques/concepts/euler-mascheroni.md} contain the definitions, derivations, and NIST DLMF citations.
\end{itemize}
\section{Learning goals}
\begin{enumerate}
  \item Define $J_n$, write $J_n=n J_{n-1}$, and compute the base $J_0=1$.
  \item Unroll to $n!$. Do not start the reduction when a substitution removes the power in one step.
  \item Convert $\int_0^1(\ln x)^n\,dx=(-1)^n n!$ by $x=e^{-t}$.
\end{enumerate}

\section{The reduction}
For an integer $n\ge 0$,
\[
J_n=\int_0^{\infty}x^n e^{-x}\,dx.
\]
The base is $J_0=\int_0^{\infty}e^{-x}\,dx=1$. For $n\ge 1$, parts with $u=x^n$ and $dv=e^{-x}\,dx$ has boundary term 0 at both ends, and
\[
J_n=n J_{n-1}.
\]
Unroll: $J_n=n!$. The index drops by one each time. Stop at 0.

\section{Worked unroll}
$J_1=1$. \texttt{U14-D10} is $J_2=2\cdot 1=2$. \texttt{U14-D02} is $J_3=3\cdot 2=6$. \texttt{U14-D09} is $J_4=4\cdot 6=24$. \texttt{U14-D01} is $J_5=5\cdot 24=120$.

\texttt{U14-D05} expands $(x+1)^2 e^{-x}=x^2 e^{-x}+2x e^{-x}+e^{-x}$, so the value is $J_2+2J_1+J_0=2+2+1=5$. \texttt{U14-D06} is the same $J_2$ reflected to $(-\infty,0]$ by $x\mapsto -x$, value 2.

\section{Faded --- a substitution, not a reduction}
\texttt{U14-D08} is $\int_0^{\infty}x^5 e^{-x^2}\,dx$. Let $u=x^2$. Then $x^5\,dx=\frac12 u^2\,du$, and the integral is $\frac12 J_2=1$. The power of the exponential changed, so the plain factorial reduction does not apply until after the substitution. \texttt{U14-D07} is the same idea with $u=\sqrt x$.

\section{Logarithms on $[0,1]$}
Let $x=e^{-t}$. Then $\int_0^1(\ln x)^n\,dx=(-1)^n J_n=(-1)^n n!$.
\texttt{U14-D04} is $n=2$, value $2$. \texttt{U14-D03} is $n=3$, value $-6$. \texttt{U14-D11} is $n=4$, value $24$. The sign is the only extra step. The reduction is the one already used for $J_n$.

\section{Speed drill --- three multiplications}
Sixty seconds. From $J_0=1$,
\[
J_2=2,\quad J_3=6,\quad J_4=24,\quad J_5=120.
\]
Those four values are \texttt{U14-D10}, \texttt{U14-D02}, \texttt{U14-D09}, and \texttt{U14-D01}. The same minute carries the sign on $[0,1]$: \texttt{U14-D04} is $+2$, \texttt{U14-D03} is $-6$, and \texttt{U14-D11} is $+24$. \texttt{U14-D05} is the sum $J_2+2J_1+J_0=2+2+1=5$. The faded blank is the four factorials, then the three signs.

\section{Optional recognition drill --- Beta}
Prerequisites: substitution and the Gamma reduction above. For real $a,b>0$,
\[
B(a,b)=\int_0^1 t^{a-1}(1-t)^{b-1}\,dt
=\frac{\Gamma(a)\Gamma(b)}{\Gamma(a+b)}.
\]
The endpoint restrictions matter. The card \path{docs/techniques/concepts/beta.md} derives the identity by $u=rt$, $v=r(1-t)$ in the product of two Gamma integrals; its Jacobian is $r$. Source: \url{https://dlmf.nist.gov/5.12.E1}.

\textbf{Worked.} The original explanatory integral
$\int_0^1 t(1-t)^2\,dt=B(2,3)=1!\,2!/4!=1/12$ also follows by expanding: $1/2-2/3+1/4=1/12$. SymPy checks both computations. It is not a contest answer.

\textbf{Faded.} Hide the worked line. From $a-1=1$ and $b-1=2$, fill the arguments in $\Gamma(\square)\Gamma(\square)/\Gamma(\square)$, then convert to factorials. The key is $2,3,5$, then $1!,2!,4!$.

\textbf{Bare.} Name the Beta integral for $t(1-t)^2$ on $[0,1]$, give its value and endpoint conditions, then decide whether $\int_0^1 t^{-1}(1-t)^2\,dt$ is finite. Accept $B(2,3)=1/12$ with $a,b>0$; refuse the variant because its first argument is 0 and it behaves like $1/t$ at zero. On an index miss, repeat the faded step; on an endpoint miss, compare with $t^p$, integrable at zero for $p>-1$.

\section{Optional recognition drill --- Euler--Mascheroni}
Prerequisites: the Gamma integral and finite sums. Let $H_n=\sum_{k=1}^n1/k$. The constant is $\gamma=\lim_{n\to\infty}(H_n-\ln n)$. Rectangle comparison shows $H_n-\ln n>0$ and
\[
(H_n-\ln n)-(H_{n+1}-\ln(n+1))
=\int_n^{n+1}\frac{dx}{x}-\frac1{n+1}>0.
\]
The sequence decreases to its limit. Source: \url{https://dlmf.nist.gov/5.2.E3}.

\textbf{Worked.} Distinguish the constant $\gamma$ from the function $\Gamma(s)$:
\[
\Gamma(1)=1,\qquad
\Gamma'(1)=\int_0^\infty e^{-x}\ln x\,dx=-\gamma.
\]
The sign uses $\Gamma'(1)=-\gamma$ from \url{https://dlmf.nist.gov/5.4.E11}. The card \path{docs/techniques/concepts/euler-mascheroni.md} supplies the differentiation argument. For $1/2\le s\le3/2$, the derivative integrand is bounded in absolute value by $x^{-1/2}|\ln x|$ on $(0,1)$ and by $e^{-x}x^{1/2}\ln x$ on $[1,\infty)$; both are integrable. SymPy checks the integral and the Gamma derivative.

\textbf{Faded.} Fill $\partial_s x^{s-1}=x^{s-1}\square$ and $\int_0^\infty e^{-x}\ln x\,dx=\square\gamma$. The blanks are $\ln x$ and $-1$. The cited special value supplies the sign; the differentiation rule alone does not.

\textbf{Bare.} Define $\gamma$, give $\Gamma(1)$ and $\Gamma'(1)$, and explain why $3/2-\ln2$ is not $\gamma$. Accept the limit, $1$, $-\gamma$, and the strict decrease from $H_2-\ln2$. On a sign miss, return to the faded step; on a limit miss, return to the rectangle comparison. Keep the justified named constant as the exact value.

Both optional recognition drills are oral checks of the exposition. They have no machine ids, are not submitted to \texttt{practice check}, and do not supply automatic tick evidence. The stored Gamma drills below assess the prerequisite, not these two optional nodes.

\section{Drill pointers}
See \texttt{drills.tex} / \texttt{drills.md}. Feed: \texttt{verify/drills\_unit14.json}. The speed drill above checks stored rows. It adds no id.

\section*{On a miss}
Teach this file. Write $J_n=n J_{n-1}$ and the base $J_0=1$, then unroll. If the exponential is $e^{-x^2}$ or $e^{-\sqrt x}$, substitute before you quote a factorial. Do not invent an answer that is not in \texttt{drills.md}.
\end{document}
